\section{Proofs for the planar walk}\label{app:planar}

This appendix proves the bound on $o^{(4)}_N$ in Proposition~\ref{prop:decomp},
Theorems~\ref{thm:planar-first} and~\ref{thm:fourdir}, and Proposition~\ref{prop:pureturns}. We
identify $\R^2$ with $\mathbb C$ and write $\mathrm i$ for the imaginary unit, so that
$e_z=\mathrm i^z$ for $z\in\Z_4$. We use the notation of Lemma~\ref{lem:threelevel}. So a path is
given by its number $m$ of switches, its skeleton $Z$ and its run lengths $\ell$, and
$Y_N=\sum_{t=0}^m\ell_t\mathrm i^{Z_t}$. Let $\mathrm{bin}(m)=\binom{N-1}m(\beta h)^mp_0^{N-1-m}$ be
the law of $m$, whose mean is $\bar m=(N-1)\beta h$. We write $\Prob_z$ and $\E_z$ for the skeleton
walk of Lemma~\ref{lem:threelevel}(c) started at $z$, $\Prob$ for its uniform start, and $\xi$ for
one of its increments. Put
$\varkappa=s/\beta$, so $\varkappa<1$ when $r>0$. For a skeleton, let $a,b,c,d$ be the numbers of
$t$ with $Z_t=0,2,1,3$, that is, the numbers of runs going east, west, north and south. So
$a+b+c+d=m+1$, and the path uses all 4 directions exactly when $a,b,c,d\ge1$. Throughout, $N=2\nu$
is even, and $[\,\cdot\,]$ is $1$ if the statement inside holds and $0$ otherwise. We also use
$\psi(z)=z(I_0(z)+I_1(z))$ and the form \eqref{eq:planar-bessel} of \cite[Thm.~5.2]{paperA}.

\subsection{Return counts and skeletons}

For $k\ge1$ and an integer $X\ge0$ let $\operatorname{comp}_k(X)=\binom{X-1}{k-1}[X\ge1]$, and let
$\operatorname{comp}_0(X)=[X=0]$. Then $\operatorname{comp}_k(X)$ is the number of compositions of
$X$ into $k$ positive parts.

\begin{lemma}[return counts]\label{lem:returncount}
Let $Z$ be a skeleton with $m$ switches and counts $a,b,c,d$. The number of compositions $\ell$ of
$N$ into $m+1$ parts with $\sum_t\ell_t\mathrm i^{Z_t}=0$ is
\[
 K_N(a,b,c,d)=\sum_{X=0}^{\nu}\operatorname{comp}_a(X)\operatorname{comp}_b(X)
 \operatorname{comp}_c(\nu-X)\operatorname{comp}_d(\nu-X).
\]
If $a,b,c,d\ge1$, then $m\ge3$ and $K_N(a,b,c,d)\le\frac12\binom N{m-2}$.
\end{lemma}

\begin{proof}
The sum $\sum_t\ell_t\mathrm i^{Z_t}$ vanishes exactly when the runs going east have the same total
length $X$ as the runs going west, and the runs going north have the same total length $X'$ as the
runs going south. Then $2X+2X'=N$, so $X'=\nu-X$. The lengths of the $a$ runs going east, in their
order, form a composition of $X$ into $a$ parts, and similarly for the other 3 directions. These 4
compositions determine $\ell$ and can be chosen independently. This gives the formula.

Now let $a,b,c,d\ge1$, so $m+1=a+b+c+d\ge4$. If $m=3$, there is 1 run in each direction, and
$K_N$ counts the pairs $X,X'\ge1$ with $X+X'=\nu$. So $K_N=\nu-1<\frac12\binom N1$. Let $m\ge4$.
Pick positions $t_0,t_1,t_2,t_3$ with $Z_{t_z}=z$ and fix the other $m-3$ run lengths. If these add
up to $N-j$, then $j\ge4$, and there are $\binom{N-j-1}{m-4}$ ways to choose them. The conditions
for $Y_N=0$ now read $\ell_{t_0}-\ell_{t_2}=\delta_1$ and $\ell_{t_1}-\ell_{t_3}=\delta_2$, where
$\delta_1$ and $\delta_2$ depend only on the fixed lengths, and the 4 free lengths add up to $j$. So
$\ell_{t_2}$ determines the other 3. Moreover $\ell_{t_0}+\ell_{t_2}=2\ell_{t_2}+\delta_1$ lies in
$[2,j-2]$ and has the parity of $\delta_1$. So there are at most $j/2$ solutions, and
\[
 K_N(a,b,c,d)\le\sum_{j=0}^{N-1}\frac j2\binom{N-1-j}{m-4}=\frac12\binom N{m-2}.
\]
The identity counts the subsets of $\{0,\dots,N-1\}$ of size $m-2$ by their second smallest element
$j$.
\end{proof}

By Lemma~\ref{lem:threelevel}, given $m$ and $Z$ the run lengths are uniform on the $\binom{N-1}m$
compositions. So $\varpi_N=K_N(a,b,c,d)/\binom{N-1}m$ is the probability that $Y_N=0$ given $m$
and $Z$, and
\begin{equation}\label{eq:app-planar-o4}
 o^{(4)}_N=\sum_{m\ge3}\mathrm{bin}(m)\,\E\bigl[\varpi_N\mathbf 1\{a,b,c,d\ge1\}\bigm|m\bigr].
\end{equation}

\begin{lemma}[skeleton returns]\label{lem:skelreturn}
Let $W_m=\sum_{t=0}^m\mathrm i^{Z_t}$, so $W_m=(a-b)+\mathrm i(c-d)$.
\begin{enumerate}[(a)]
\item $\E_z\lvert W_m\rvert^2\le m+1$ for every start $z$.
\item If $s>0$ and $m$ is odd, then $\Prob(W_m=0)\ge\varkappa/(36(m+1))$.
\item If $s=0$ and $m\equiv3\pmod4$, then $\Prob(W_m=0)\ge1/(m+1)$.
\end{enumerate}
\end{lemma}

\begin{proof}
(a) For $t>t'$ the difference $Z_t-Z_{t'}$ is a sum of $t-t'$ independent copies of $\xi$, and
$\E\,\mathrm i^\xi=\frac r\beta(\mathrm i+\mathrm i^{-1})+\frac s\beta\mathrm i^2=-\varkappa$. So
\[
 \E_z\lvert W_m\rvert^2=\sum_{t,t'=0}^m\E_z\,\mathrm i^{Z_t-Z_{t'}}
 =m+1+2\sum_{k=1}^m(m+1-k)(-\varkappa)^k .
\]
The last sum has terms of alternating sign, the first one negative, and their absolute values do
not increase since $0\le\varkappa\le1$. So it is at most $0$.

(b) Write $m=2m_1+1$ and split $W_m=W'+W''$, where $W'$ sums over $t\le m_1$ and $W''$ over
$t>m_1$. The law of $\xi$ is symmetric on $\Z_4$ and the uniform law is stationary, so
$(Z_{m_1},Z_{m_1-1},\dots,Z_0)$ is again the walk from a uniform start. Hence
$\Prob(W'=w,\,Z_{m_1}=z)=\frac14\Prob_z(W_{m_1}=w)$. The event $Z_{m_1+1}=Z_{m_1}+2$ has probability
$\varkappa$ and is independent of $Z_0,\dots,Z_{m_1}$. On this event, given $Z_{m_1}=z$, the second
half is the walk from $z+2$, and adding $2$ to every entry changes the sign of $W''$. So
$\Prob_{z+2}(W_{m_1}=-w)=\Prob_z(W_{m_1}=w)$ and
\[
 \Prob(W_m=0)\ge\frac\varkappa4\sum_{z\in\Z_4}\sum_w\Prob_z(W_{m_1}=w)^2 .
\]
Let $R^2=m+1=2(m_1+1)$. By (a) and Chebyshev's inequality, $\Prob_z(\lvert W_{m_1}\rvert\le R)\ge\frac12$.
At most $(2R+1)^2\le9R^2$ Gaussian integers $w$ have $\lvert w\rvert\le R$, so by the
Cauchy--Schwarz inequality the inner sum is at least $\frac14/(9R^2)$. Summing over the 4 values
of $z$ gives $\Prob(W_m=0)\ge\varkappa/(36(m+1))$.

(c) Now every switch is a turn, so $\mathrm i^{Z_t}=\mathrm i^{Z_{t-1}}\cdot\mathrm i\epsilon_t$
with independent fair signs $\epsilon_t$. Hence $\mathrm i^{Z_t}=\mathrm i^{Z_0}\mathrm i^t\sigma_t$,
where $\sigma_0=1$ and $\sigma_t=\epsilon_1\cdots\epsilon_t$, and $\sigma_1,\dots,\sigma_m$ are
independent fair signs. Write $m=4k-1$. Then $\mathrm i^{-Z_0}W_m$ has real part
$\sum_{t\text{ even}}(-1)^{t/2}\sigma_t$, which is $1$ plus a sum of $2k-1$ independent fair signs,
and imaginary part $\sum_{t\text{ odd}}(-1)^{(t-1)/2}\sigma_t$, an independent sum of $2k$ fair
signs. So
\[
 \Prob(W_m=0)=\binom{2k-1}k2^{-(2k-1)}\binom{2k}k2^{-2k}=\Bigl(\binom{2k}k4^{-k}\Bigr)^2
 =\prod_{j=1}^k\frac{(2j-1)^2}{(2j)^2}\ge\frac14\prod_{j=2}^k\frac{2j-2}{2j}=\frac1{4k},
\]
using $(2j-1)^2\ge(2j-2)2j$. This is $1/(m+1)$.
\end{proof}

\begin{lemma}[balanced skeletons]\label{lem:balanced}
If $a=b\ge1$, $c=d\ge1$ and $m=2(a+c)-1\le N-1$, then
$K_N(a,b,c,d)/\binom{N-1}m\ge\sqrt{m+1}/(2N^2)$.
\end{lemma}

\begin{proof}
Put $k=a+c=(m+1)/2$, so $2\le k\le\nu$. Only $1\le X\le\nu-1$ contribute to $K_N$. So by the
Cauchy--Schwarz inequality and Vandermonde's identity,
\[
 K_N=\sum_{X=1}^{\nu-1}\binom{X-1}{a-1}^2\binom{\nu-X-1}{c-1}^2
 \ge\frac1\nu\Bigl(\sum_{X=1}^{\nu-1}\binom{X-1}{a-1}\binom{\nu-X-1}{c-1}\Bigr)^2
 =\frac1\nu\binom{\nu-1}{k-1}^2 .
\]
Next, $\binom{\nu-1}{k-1}^2=\frac k\nu\binom{\nu-1}{k-1}\binom\nu k$, and
$\binom{\nu-1}{k-1}\binom\nu k/\binom{2\nu-1}{2k-1}=\Prob(H=k-1)$, where $H$ is the number of red
balls among $2k-1$ balls drawn without replacement from $\nu-1$ red and $\nu$ blue ones. If $k=\nu$,
all balls are drawn and $H=k-1$. Otherwise the ratio
\[
 \frac{\Prob(H=j+1)}{\Prob(H=j)}=\frac{(\nu-1-j)(2k-1-j)}{(j+1)(\nu-2k+2+j)}
\]
decreases in $j$, equals $\frac{\nu-k}{\nu-k+1}<1$ at $j=k-1$ and exceeds $1$ at $j=k-2$. So $k-1$
is a mode of $H$. Since $\Var H\le(2k-1)/4$, Chebyshev's inequality puts mass at least $\frac34$ on
the at most $2\sqrt{2k-1}+1\le3\sqrt{2k}$ integers within $\sqrt{2k-1}$ of $\E H$. So
$\Prob(H=k-1)\ge1/(4\sqrt{2k})$. As $\binom{N-1}m=\binom{2\nu-1}{2k-1}$ and $\nu=N/2$,
\[
 \frac{K_N}{\binom{N-1}m}\ge\frac k{\nu^2}\cdot\frac1{4\sqrt{2k}}=\frac{\sqrt{m+1}}{2N^2}.
 \qedhere
\]
\end{proof}

\subsection{The four-direction bound and the first order}

\begin{proof}[Proof of the bound on $o^{(4)}_N$ in Proposition~\ref{prop:decomp}]
Note that $(\beta h)^mp_0^{N-1-m}=p_0^{N-1}w^m$. By \eqref{eq:app-planar-o4} and
Lemma~\ref{lem:returncount},
\[
 o^{(4)}_N\le\sum_{m=3}^{N-1}\binom{N-1}m(\beta h)^mp_0^{N-1-m}\,\frac{\binom N{m-2}}{2\binom{N-1}m}
 =\frac{p_0^{N-1}}2\sum_{m=3}^{N-1}\binom N{m-2}w^m
 \le\frac{p_0^{N-1}w^2}2\bigl((1+w)^N-1\bigr).
\]
Since $1+w=1/p_0$, the right side equals $\frac{w^2}{2p_0}(1-p_0^N)$. Dividing by
$c_N=\frac14p_0^{N-1}$ gives $R_N=o^{(4)}_N/c_N\le2w^2((1+w)^N-1)$.
\end{proof}

\begin{proof}[Proof of Theorem~\ref{thm:planar-first}]
We first prove 3 bounds. Since $I_1\le I_0\le e^z$ for $z\ge0$, we have $\psi(z)\le2ze^z$, and
\eqref{eq:planar-bessel} gives $2S_N(u)\le8ue^{Nu}$. With Proposition~\ref{prop:decomp} and
$(1+w)^N\le e^{Nw}$,
\begin{equation}\label{eq:app-planar-upper}
 \frac{o_N}{c_N}\le8ue^{Nu}+2w^2e^{Nw}.
\end{equation}
Since $o_N\ge o^{(0)}_N$, \eqref{eq:planar-bessel} with $z=Nu$ gives
\begin{equation}\label{eq:app-planar-axes}
 \frac{o_N}{c_N}\ge2S_N(u)\ge\frac4N\Bigl(1-\frac{z(z+2)}{2N}\Bigr)\psi(z).
\end{equation}
For the third bound we show that there are $c'>0$ and $T_0$, depending only on $r$ and $s$, with
\begin{equation}\label{eq:app-planar-four}
 o^{(4)}_N\ge\frac{c'}{N^2\sqrt T}\qquad\text{whenever }T=Nh\ge T_0\text{ and }\beta h\le\tfrac12 .
\end{equation}
Let $\mathcal M$ be the set of odd $m$ if $s>0$, and the set of $m\equiv3\pmod4$ if $s=0$. Let
$m\in\mathcal M$. The event $\{a=b\ge1,\ c=d\ge1\}$ is $\{W_m=0\}$ minus the events $\{a=b=0\}$ and
$\{c=d=0\}$. Each of these needs every switch to be a reversal, so it has probability at most
$\frac12\varkappa^m$, and it is empty if $s=0$. If $s>0$, choose $m_2$ with
$\varkappa^m\le\varkappa/(72(m+1))$ for $m\ge m_2$, and put $c_2=\varkappa/72$. If $s=0$, put
$m_2=3$ and $c_2=1$. By Lemma~\ref{lem:skelreturn}, $\Prob(a=b\ge1,\ c=d\ge1\mid m)\ge c_2/(m+1)$
for $m\in\mathcal M$ with $m\ge m_2$. On this event Lemma~\ref{lem:balanced} gives
$\varpi_N\ge\sqrt{m+1}/(2N^2)$. Let $J=[\bar m/2,2\bar m]$. By \eqref{eq:app-planar-o4}, if
$\bar m\ge2m_2$ then
\[
 o^{(4)}_N\ge\sum_{m\in\mathcal M\cap J}\mathrm{bin}(m)\,\frac{c_2}{2N^2\sqrt{m+1}}
 \ge\frac{c_2}{2N^2\sqrt{2\bar m+1}}\,\Prob(m\in\mathcal M\cap J).
\]
By the Chernoff bound, $\Prob(m\notin J)\le e^{-\bar m/8}+e^{-\bar m/3}\le2e^{-\bar m/8}$. Put
$p=\beta h\le\frac12$. Then $\Prob(m\text{ odd})=\frac12(1-(1-2p)^{N-1})\ge\frac12(1-e^{-2\bar m})$.
For $s=0$ we use the roots of unity filter,
\[
 \Prob(m\equiv3\bmod4)=\frac14\sum_{k=0}^3\mathrm i^{-3k}\bigl(1-p+p\,\mathrm i^k\bigr)^{N-1}
 \ge\frac14-\frac34e^{-\bar m/2},
\]
since $\lvert1-p+p\,\mathrm i^{\pm1}\rvert^2=1-2p(1-p)\le e^{-p}$ and $\lvert1-2p\rvert\le e^{-2p}$. So in
both cases $\Prob(m\in\mathcal M\cap J)\ge\frac14-\frac34e^{-\bar m/2}-2e^{-\bar m/8}$, which is at
least $\frac18$ when $\bar m$ is large. Since $\bar m\ge\beta T/2$, and $2\bar m+1\le3\beta T$ when
$\beta T\ge1$, this proves \eqref{eq:app-planar-four} with $c'=c_2/(16\sqrt{3\beta})$.

(a) Let $h=\tau L/N$ and $T=\tau L$. Then $u,w=O(L/N)$, $Nu=sT/p_0=s\tau L+O(L^2/N)$ and
$Nw=\beta T/p_0=\beta\tau L+O(L^2/N)$. So \eqref{eq:app-planar-upper} gives
$o_N/c_N=O(LN^{s\tau-1}+L^2N^{\beta\tau-2})$. If $\tau<\tau^*$, then $s\tau<1$ and $\beta\tau<2$,
and this is at most $N^{-\gamma}$ for large $N$ whenever $0<\gamma<\min(1-s\tau,2-\beta\tau)$. Now
let $\tau>\tau^*$, so that $s\tau>1$ or $\beta\tau>2$. If $s\tau>1$, then $s>0$, $z=Nu\to\infty$ and
$z^2/N\to0$, and \eqref{eq:planar-psi} gives $\psi(z)=\sqrt{2z/\pi}\,e^z(1+o(1))$. So
\eqref{eq:app-planar-axes} gives $o_N/c_N\ge(1+o(1))\frac4N\sqrt{2z/\pi}\,e^z$, which is at least
$N^{s\tau-1}$ for large $N$. If $\beta\tau>2$, we use \eqref{eq:app-planar-four} and
$c_N^{-1}=4p_0^{-(N-1)}\ge4e^{\beta h(N-1)}=4e^{\beta T-\beta h}$. They give
$o_N/c_N\ge4c'e^{-\beta h}N^{\beta\tau-2}(\tau L)^{-1/2}$, which is at least $N^{(\beta\tau-2)/2}$
for large $N$.

(b) Let $\tau<\tau^*<\tau'$. By (a), for large $N$ the ratio $o_N/c_N$ is less than $1$ at
$h=\tau L/N$ and greater than $1$ at $h=\tau'L/N<1/\beta$. By Proposition~\ref{prop:decomp} it
increases in $h$, so $\tau L<T^*_N<\tau'L$.

(c) Put $\tau_N=T^*_N/L$, so $\tau_N\to\tau^*$ by (b), and let $u^*=u^*_N$ and
$w^*=\beta h^*_N/(1-\beta h^*_N)$. As in (a), $u^*,w^*=O(L/N)$, $Nu^*=s\tau_NL+o(1)$ and
$Nw^*=\beta\tau_NL+o(1)$. At the crossing Proposition~\ref{prop:decomp} gives $2S_N(u^*)+R_N=1$.
If $s>2r$, then $\tau^*=1/s$ and $\beta\tau^*=\beta/s<2$. So
$R_N\le2w^{*2}e^{Nw^*}=N^{\beta/s-2+o(1)}\to0$, and $2S_N(u^*)\to1$. The 2 axes carry the same mass
by symmetry. If $s<2r$, then $\tau^*=2/\beta$ and $s\tau^*=2s/\beta<1$. So
$2S_N(u^*)\le8u^*e^{Nu^*}=N^{2s/\beta-1+o(1)}$, which is at most $N^{-\gamma}$ for large $N$
whenever $0<\gamma<1-2s/\beta$.
\end{proof}

\subsection{The four-direction local limit}

Given the skeleton, the run lengths are close to the gaps between $m$ uniform points in $[0,N]$.
So the total lengths in the 4 directions are close to $N$ times a Dirichlet vector
$(y_0,y_1,y_2,y_3)$ with parameters $(a,c,b,d)$. Then $N^2\varpi_N$ should be close to twice the
density of $(y_0-y_2,y_1-y_3)$ at $0$, where the factor $2$ is the parity of $Y_N$. This is the
quantity $g(a,b,c,d)$ below. Divided by $m$, it is close to a Gaussian function of $W_m/\sqrt m$,
and $W_m$ satisfies a central limit theorem.

\begin{proof}[Proof of Theorem~\ref{thm:fourdir}]
Let $A=a+b$ and $B=c+d$, so $A+B=m+1$, and put $\Gamma_4=(a-1)!\,(b-1)!\,(c-1)!\,(d-1)!$ and
\[
 g(a,b,c,d)=\frac{m(m-1)}2\cdot\frac{\binom{A-2}{a-1}}{2^{A-2}}\cdot\frac{\binom{B-2}{c-1}}{2^{B-2}}
 =\frac{m!\,(A-2)!\,(B-2)!}{2^{m-2}\,(m-2)!\,\Gamma_4}.
\]
Let $J_N=[\bar m/2,2\bar m]$, $\mathcal G_m=\{a,b,c,d\ge m/8\}$ and $m_0=20$. Note that
$\bar m=\beta T(1-1/N)\to\infty$ and $\bar m\le\beta N^{1/8}$. Below, $C$ is a constant that depends
only on $r$ and $s$ and may change from line to line.

\emph{Step 1 (few or many switches).} If $a,b,c,d\ge1$ and $m\le N/5$, then
Lemma~\ref{lem:returncount} and
$\binom N{m-2}/\binom{N-1}m=Nm(m-1)/((N-m+2)(N-m+1)(N-m))$ give $\varpi_N\le m^2/N^2$. If $m>N/5$,
then $\varpi_N\le1<25m^2/N^2$. So $N^2\varpi_N\le25m^2$ whenever the path uses all 4 directions.
By the Chernoff bound, $\Prob(m\notin J_N)\le2e^{-\bar m/8}$. Also
$\E m^4=\sum_{j=1}^4S(4,j)(N-1)_j(\beta h)^j\le15\bar m^4$, where $S(4,j)=1,7,6,1$ are Stirling
numbers of the second kind, $(N-1)_j$ is a falling factorial, and $\bar m\ge1$. By the
Cauchy--Schwarz inequality, the terms of \eqref{eq:app-planar-o4} with $m\notin J_N$ contribute at
most $25\,\E[m^2\mathbf 1\{m\notin J_N\}]\le25\sqrt{30}\,\bar m^2e^{-\bar m/16}=o(1)$ to
$N^2o^{(4)}_N$.

\emph{Step 2 (unbalanced skeletons).} For $k=1,2,3$ let $V_k=\sum_{t=0}^m\mathrm i^{kZ_t}$. The
number of $t$ with $Z_t=z$ is $\frac{m+1}4+\frac14\sum_{k=1}^3\mathrm i^{-kz}V_k$. Let
$\vartheta_k=\E\,\mathrm i^{k\xi}$, so $\vartheta_1=\vartheta_3=-\varkappa$ and
$\vartheta_2=2\varkappa-1$. Then
$\E[\mathrm i^{kZ_t}\mid Z_0,\dots,Z_{t-1}]=\vartheta_k\mathrm i^{kZ_{t-1}}$, so
$M^{(k)}_m=\sum_{t=1}^m(\mathrm i^{kZ_t}-\vartheta_k\mathrm i^{kZ_{t-1}})$ is a martingale with
increments of modulus at most $2$, and
\[
 (1-\vartheta_k)V_k=M^{(k)}_m+\mathrm i^{kZ_0}-\vartheta_k\mathrm i^{kZ_m}.
\]
Here $1-\vartheta_k\ge\min(1,4r/\beta)>0$. Burkholder's inequality \cite[Ch.~2]{hallheyde1980},
applied to the real and imaginary parts, gives $\E\lvert M^{(k)}_m\rvert^4\le Cm^2$. So
$\E\lvert V_k\rvert^4\le Cm^2$ and $\E(a-\frac{m+1}4)^4\le Cm^2$, and the same holds for $b,c,d$.
Since $\frac{m+1}4-\frac m8\ge\frac m8$, Markov's inequality gives
$\Prob(\mathcal G_m^c\mid m)\le C/m^2$. Every skeleton in $\mathcal G_m$ uses all 4 directions. By
Step 1, the skeletons with all 4 directions that are not in $\mathcal G_m$, with $m\in J_N$,
contribute at most $\sum_m\mathrm{bin}(m)\,25m^2\,C/m^2\le25C$ to $N^2o^{(4)}_N$.

\emph{Step 3 (run lengths given the skeleton).} We show that $N^2\varpi_N=g(a,b,c,d)(1+o(1))$
uniformly over the skeletons in $\mathcal G_m$ with $m\in J_N$. For large $N$ such $m$ satisfy
$m_0\le m\le2\beta N^{1/8}\le N^{1/6}$. On $\mathcal G_m$ we have $A,B\ge m/4$, so
$A-1,B-1\ge m/8$. Put $f(X)=X^{A-2}(\nu-X)^{B-2}$. For
$X\ge1$ we have $\operatorname{comp}_a(X)=\frac{X^{a-1}}{(a-1)!}\prod_{j<a}(1-j/X)_+$, and
similarly for $b,c,d$. So
\[
 K_N(a,b,c,d)=\frac1{\Gamma_4}\sum_{X=1}^{\nu-1}f(X)\chi(X),
\]
where $\chi(X)\in[0,1]$ is the product of the 4 products. The function $f$ vanishes at $0$ and
$\nu$ and is unimodal, and its mode $X^*=\nu(A-2)/(m-3)$ lies in $[\nu/6,5\nu/6]$ since
$A,B\ge m/4$ and $m\ge18$. On each monotone piece of $f$ every value $f(X)$ lies between the integrals of $f$ over the 2 unit
intervals next to $X$. So $\lvert\sum_{j\le X\le k}f(X)-\int_j^kf\rvert\le2\max f$ for all integers
$0\le j\le k\le\nu$. Moreover $f/\int_0^\nu f$ is the density of $\nu$ times a Beta variable with
parameters $A-1$ and $B-1$, and $\int_0^\nu f=\nu^{m-2}\mathrm B(A-1,B-1)$.

We need 3 estimates. First, this Beta law is the law of the $(A-1)$-th smallest of $m-2$
independent uniform points in $[0,1]$. A union bound over sets of $A-1$ points gives, for $0<t<1$,
\[
 \Prob(\mathrm{Beta}\le t)\le\binom{m-2}{A-1}t^{A-1}\le\Bigl(\frac{e(m-2)t}{A-1}\Bigr)^{A-1}
 \le(8et)^{A-1}.
\]
We take $t=2\nu^{-1/4}$. Since $A-1\ge2$ and $16e\nu^{-1/4}<1$ for large $N$, this gives
$\int_0^{\nu^{3/4}+1}f\le(16e)^2\nu^{-1/2}\int_0^\nu f$, and the same bound holds near $\nu$. Second,
for $\nu^{3/4}\le X\le\nu-\nu^{3/4}$, the inequality $1-\prod_j(1-t_j)_+\le\sum_jt_j$ for $t_j\ge0$
gives $1-\chi(X)\le(a^2+b^2+c^2+d^2)/(2\nu^{3/4})\le m^2\nu^{-3/4}$. Third, on
$[\nu/12,11\nu/12]$ we have
$(\log f)''=-(A-2)/X^2-(B-2)/(\nu-X)^2\ge-144m/\nu^2$. This interval contains every $X$ with
$\lvert X-X^*\rvert\le\nu/(12\sqrt m)$, and for these $X$ Taylor's formula at $X^*$ gives
$f(X)\ge e^{-1/2}\max f$. So $\int_0^\nu f\ge e^{-1/2}\max f\cdot\nu/(6\sqrt m)$, that is,
$\max f\le6e^{1/2}\sqrt m\,\nu^{-1}\int_0^\nu f$.

Now $\sum_{X=1}^{\nu-1}f\chi\le\sum_Xf\le\int_0^\nu f+2\max f$. Summing only over
$\nu^{3/4}\le X\le\nu-\nu^{3/4}$ gives
\[
 \sum_{X=1}^{\nu-1}f\chi\ge\bigl(1-m^2\nu^{-3/4}\bigr)
 \Bigl(\int_0^\nu f-2(16e)^2\nu^{-1/2}\int_0^\nu f-2\max f\Bigr).
\]
Since $m^2\le N^{1/3}$ and $\sqrt m/\nu\le N^{-5/6}$, we get
$\Gamma_4K_N=\nu^{m-2}\mathrm B(A-1,B-1)(1+o(1))$, uniformly. Also
$\binom{N-1}m=\frac{N^m}{m!}\prod_{j=1}^m(1-\frac jN)=\frac{N^m}{m!}(1+O(m^2/N))$. With
$\nu^{m-2}/N^{m-2}=2^{-(m-2)}$ and $\mathrm B(A-1,B-1)=(A-2)!\,(B-2)!/(m-2)!$ this gives
$N^2\varpi_N=g(a,b,c,d)(1+o(1))$.

\emph{Step 4 (average over the skeleton).} We show that
$\E[g\,\mathbf 1_{\mathcal G_m}\mid m]/m\to c_*=\frac{2(r+s)}{\pi\beta}$ as $m\to\infty$. First we
need a local limit for the fair binomial law. Let $0\le k\le n$ and $y=(2k-n)/n$. If
$\lvert2k-n\rvert\le n^{3/5}$, Stirling's formula gives
\begin{align*}
 \sqrt n\binom nk2^{-n}&=\sqrt{\frac2{\pi(1-y^2)}}\,
 e^{-\frac n2[(1+y)\log(1+y)+(1-y)\log(1-y)]}\bigl(1+O(\tfrac1n)\bigr)\\
 &=\sqrt{\frac2\pi}\,e^{-\frac{(2k-n)^2}{2n}}\bigl(1+O(n^{-3/5})\bigr),
\end{align*}
since the bracket is $y^2+O(y^4)$ and $ny^4\le n^{-3/5}$. Both sides decrease as $\lvert2k-n\rvert$
grows, so outside this range both are $O(e^{-n^{1/5}/3})$. Hence
\begin{equation}\label{eq:app-planar-dml}
 \sup_{0\le k\le n}\Bigl\lvert\sqrt n\binom nk2^{-n}-\sqrt{\tfrac2\pi}\,e^{-(2k-n)^2/(2n)}\Bigr\rvert
 =O(n^{-1/2}).
\end{equation}
On $\mathcal G_m$ we have $A-2,B-2\ge m/8$. We apply \eqref{eq:app-planar-dml} with $n=A-2$,
$k=a-1$ and with $n=B-2$, $k=c-1$, and note that $2(a-1)-(A-2)=a-b$. Then, uniformly on
$\mathcal G_m$,
\[
 \frac gm=\frac{m-1}{\pi\sqrt{(A-2)(B-2)}}
 \exp\Bigl(-\frac{(a-b)^2}{2(A-2)}-\frac{(c-d)^2}{2(B-2)}\Bigr)+O(m^{-1/2}).
\]
Now we replace $m-1$ by $m$, then $A-2$ by $A$ and $B-2$ by $B$, and finally $B$ by $m(1-\theta)$,
where $\theta=A/m$. The last step is needed because $B=m(1-\theta)+1$. Each step changes the
prefactor by a factor $1+O(1/m)$ and each exponential by $O(1/m)$, since
$0\le e^{-y/A_2}-e^{-y/A_1}\le\frac{A_2-A_1}{eA_1}$ for $y\ge0$ and $0<A_1\le A_2$. So, uniformly on
$\mathcal G_m$,
\[
 \frac gm=\Psi\Bigl(\frac{a-b}{\sqrt m},\frac{c-d}{\sqrt m},\frac Am\Bigr)+O(m^{-1/2}),\qquad
 \Psi(x_1,x_2,\theta)=\frac{e^{-x_1^2/(2\theta)-x_2^2/(2(1-\theta))}}{\pi\sqrt{\theta(1-\theta)}}.
\]
On $\mathcal G_m$ we have $\theta\in[\frac14,\frac34+\frac1m]\subset[\frac15,\frac45]$, as $m\ge20$.
Let $\tilde\Psi$ be $\Psi$ with $\theta$ replaced by the nearest point of $[\frac15,\frac45]$. It is
bounded and continuous on $\R^3$ and equals $\Psi$ on $\mathcal G_m$.

Next we prove a central limit theorem for $W_m=(a-b)+\mathrm i(c-d)$. Put
$\zeta_j=\mathrm i^{\xi_j}+\varkappa$, where $\xi_j=Z_j-Z_{j-1}$. The $\zeta_j$ are independent with
$\E\zeta_j=0$, $\E\lvert\zeta_j\rvert^2=1-\varkappa^2$ and
$\E\zeta_j^2=\vartheta_2-\varkappa^2=-(1-\varkappa)^2$. Since
$\mathrm i^{Z_j}=\mathrm i^{Z_{j-1}}(\zeta_j-\varkappa)$,
\[
 (1+\varkappa)W_m=M_m+\mathrm i^{Z_0}+\varkappa\,\mathrm i^{Z_m},\qquad
 M_m=\sum_{j=1}^m\mathrm i^{Z_{j-1}}\zeta_j,
\]
and $M_m$ is a martingale with increments of modulus at most $2$. We take the skeletons for all $m$
to be the first steps of 1 infinite walk, so the $\sigma$-fields are nested. Fix $\varphi\in\R$. For a
constant $\omega$ with $\lvert\omega\rvert=1$ we have
$\E(\operatorname{Re}\omega\zeta_j)^2=\frac12(\E\lvert\zeta_j\rvert^2+\operatorname{Re}\omega^2\E\zeta_j^2)$.
So the $j$-th increment of $\operatorname{Re}(e^{-\mathrm i\varphi}M_m)$ has conditional variance
$\frac12[(1-\varkappa^2)-(1-\varkappa)^2\cos(2\varphi)(-1)^{Z_{j-1}}]$. The average of these
variances tends to $\frac12(1-\varkappa^2)$ in probability, because
$\sum_{t=0}^m(-1)^{Z_t}=V_2$ and $\E\lvert V_2/m\rvert^4\le C/m^2$ by Step 2. The increments are
bounded, so the martingale central limit theorem \cite[Cor.~3.1]{hallheyde1980} and the
Cram\'er--Wold device give $M_m/\sqrt m\Rightarrow N(0,\frac12(1-\varkappa^2)I_2)$. Hence
\[
 \frac{W_m}{\sqrt m}\Rightarrow X\sim N(0,\sigma_*^2I_2),\qquad
 \sigma_*^2=\frac{1-\varkappa}{2(1+\varkappa)}=\frac r{2(r+s)}.
\]
Also $A=\frac{m+1}2+\frac12V_2$, so $A/m\to\frac12$ in probability, and
$(W_m/\sqrt m,A/m)\Rightarrow(X,\frac12)$. By Step 2,
$\E[\tilde\Psi\,\mathbf 1_{\mathcal G_m^c}\mid m]\le C/m^2$. So, since $\tilde\Psi$ is bounded and
continuous,
\[
 \frac{\E[g\,\mathbf 1_{\mathcal G_m}\mid m]}m
 =\E\,\tilde\Psi\Bigl(\frac{W_m}{\sqrt m},\frac Am\Bigr)+O(m^{-1/2})
 \to\E\,\Psi\bigl(X,\tfrac12\bigr)=\frac2\pi\,\E e^{-\lvert X\rvert^2}
 =\frac2{\pi(1+2\sigma_*^2)}=c_*.
\]

\emph{Step 5 (conclusion).} By Steps 1 and 2,
$N^2o^{(4)}_N=\sum_{m\in J_N}\mathrm{bin}(m)\,\E[N^2\varpi_N\mathbf 1_{\mathcal G_m}\mid m]+O(1)$.
By Step 3 the expectation is $\E[g\,\mathbf 1_{\mathcal G_m}\mid m](1+o(1))$, with $o(1)$ uniform in
$m\in J_N$. Since $\min J_N\to\infty$, Step 4 turns this into
$N^2o^{(4)}_N=c_*(1+o(1))\sum_{m\in J_N}\mathrm{bin}(m)\,m+O(1)$. As in Step 1,
$\E[m\mathbf 1\{m\notin J_N\}]\to0$, and $\bar m=\beta T(1+O(1/N))$. Hence
$N^2o^{(4)}_N=c_*\beta T(1+o(1))+O(1)=\frac{2(r+s)}\pi T(1+o(1))$.
\end{proof}

\subsection{Turns only}

\begin{proof}[Proof of Proposition~\ref{prop:pureturns}]
Let $s=0$, so $\beta=2r$ and $p_0=1-2rh$. Step $t$ changes $U=Y^{(1)}+Y^{(2)}$ by a sign
$\sigma^U_t$ and $V=Y^{(1)}-Y^{(2)}$ by a sign $\sigma^V_t$. The directions east, north, west and
south give the sign pairs $(\sigma^U_t,\sigma^V_t)=(+,+)$, $(+,-)$, $(-,-)$ and $(-,+)$. So
directions and sign pairs correspond one to one, and $Y_N=0$ exactly when $\sum_t\sigma^U_t=\sum_t\sigma^V_t=0$. The 2 neighbors of a direction
differ from it in exactly 1 sign each, and in different signs. So at each step the walk changes
only $\sigma^U$ with probability $rh$, only $\sigma^V$ with probability $rh$, and neither with
probability $1-2rh$. Hence a path is given by 2 words in $\pm$ with first letters $\sigma,\sigma'$
and switch sets $A,B\subseteq\{1,\dots,N-1\}$ with $A\cap B=\emptyset$, and each such quadruple
gives exactly 1 path. Its probability is
$\frac14(1-2rh)^{N-1-\lvert A\rvert-\lvert B\rvert}(rh)^{\lvert A\rvert+\lvert B\rvert}
=c_Nx^{\lvert A\rvert+\lvert B\rvert}$. So $o_N/c_N$ is the sum of
$x^{\lvert A\rvert+\lvert B\rvert}$ over the pairs of balanced words with disjoint switch sets. By
the definition of $S_N$, the sum over all pairs of balanced words is $S_N(x)^2$. This proves
$o_N/c_N=S_N(x)^2-\Delta_N(x)$.

The 4 pairs with $A=B=\{\nu\}$ are balanced, since a single switch after step $\nu$ gives a
balanced word for either first letter. So $\Delta_N(x)\ge4x^2$. For the upper bound, a union bound
over a common switch $i\in A\cap B$ gives $\Delta_N(x)\le\sum_{i=1}^{N-1}S_{N,i}(x)^2$, where
$S_{N,i}(x)$ is the sum of $x^{\lvert A\rvert}$ over the balanced words $(\sigma,A)$ with $i\in A$.
If $A=\{t_1<\dots<t_k\}$, the word has sum $\sigma\bigl(2\sum_j(-1)^{j-1}t_j+(-1)^kN\bigr)$. So being
balanced does not depend on $\sigma$, and it is a linear equation in $t_1,\dots,t_k$ with
coefficients $\pm2$. For $k=1$ it gives $A=\{\nu\}$. Let $k\ge2$ and $i\in A$, and let $t'$ be the
largest element of $A\setminus\{i\}$. There are at most $\binom{N-2}{k-2}$ choices for
$A\setminus\{t',i\}$. Given this set and whether $t'>i$, the position in $A$ of every element is
known, since $t'$ is the largest or the second largest element. So the equation determines $t'$.
Hence at most $2\binom{N-2}{k-2}$ sets $A\ni i$ of size $k$ give balanced words, and
\[
 S_{N,i}(x)\le2x\,[i=\nu]+4\sum_{k\ge2}\binom{N-2}{k-2}x^k=2x\,[i=\nu]+4x^2(1+x)^{N-2}.
\]
Summing the squares gives
$\Delta_N(x)\le\bigl(2x+4x^2(1+x)^{N-2}\bigr)^2+16(N-2)x^4(1+x)^{2(N-2)}$, which is the stated upper
bound.

Now let $N\ge4$ and $x^*=x^*_N$. At the crossing, $S_N(x^*)^2=1+\Delta_N(x^*)>1=S_N(u_N)^2$, so
$x^*>u_N$. The polynomial $S_N$ has nonnegative coefficients and no constant term, so $S_N(x)/x$ is
nondecreasing, and
\[
 \frac{x^*}{u_N}\le\frac{S_N(x^*)}{S_N(u_N)}=\sqrt{1+\Delta_N(x^*)}\le1+\tfrac12\Delta_N(x^*).
\]
By Theorem~\ref{thm:planar-first}(b) with $\tau^*=1/r$, $rT^*_N=L(1+o(1))$. So
$Nx^*=rT^*_N/(1-2rT^*_N/N)=L(1+o(1))$, and with $(1+x)^n\le e^{nx}$ the upper bound on $\Delta_N$
gives
\[
 \Delta_N(x^*)\le4x^{*2}+16x^{*3}e^{Nx^*}+16Nx^{*4}e^{2Nx^*}=N^{-1+o(1)}.
\]
Hence $N(x^*-u_N)\le Nu_N\,N^{-1+o(1)}\to0$, since $Nu_N=O(L)$. Now we can be more precise. By
\cite[Eq.~(7)]{paperA}, $Nu_N=N(1-p_N)/p_N=L-\frac12\log L+c_0+o(1)$, and so also
$Nx^*=L-\frac12\log L+c_0+o(1)$. Then $e^{2Nx^*}=\frac\pi8N^2L^{-1}(1+o(1))$ and
$x^*=\frac LN(1+o(1))$. So the last term of the bound is $2\pi L^3N^{-1}(1+o(1))$, while the first 2
terms are $O(L^{5/2}N^{-2})$. This gives $\Delta_N(x^*)\le2\pi L^3N^{-1}(1+o(1))$ and the bound on
$x^*/u_N$. Finally, $rT^*_N=Nx^*/(1+2x^*)$ and $N(1-p_N)=Nu_N/(1+u_N)$, so
\[
 rT^*_N-N(1-p_N)=\frac{N(x^*-u_N)}{1+2x^*}+Nu_N\Bigl(\frac1{1+2x^*}-\frac1{1+u_N}\Bigr)
 =O\Bigl(\frac{L^4}N\Bigr)+O\Bigl(\frac{L^2}N\Bigr).
 \qedhere
\]
\end{proof}
